\documentclass{article}

\usepackage{amsmath}
\usepackage{amssymb}
\usepackage{polski}

\usepackage[utf8]{inputenc}
\usepackage[T1]{fontenc}

\usepackage{tikz}

\usepackage{listings}
\lstset{
    language=SQL,
    inputencoding=utf8x,
    extendedchars=\true,
    literate={ą}{{\k{a}}}1
             {Ą}{{\k{A}}}1
             {ę}{{\k{e}}}1
             {Ę}{{\k{E}}}1
             {ó}{{\'o}}1
             {Ó}{{\'O}}1
             {ś}{{\'s}}1
             {Ś}{{\'S}}1
             {ł}{{\l{}}}1
             {Ł}{{\L{}}}1
             {ż}{{\.z}}1
             {Ż}{{\.Z}}1
             {ź}{{\'z}}1
             {Ź}{{\'Z}}1
             {ć}{{\'c}}1
             {Ć}{{\'C}}1
             {ń}{{\'n}}1
             {Ń}{{\'N}}1
}

\newcommand\tikzmark[1]{%
\tikz[remember picture,baseline] \node[inner sep=0,outer sep=0] (#1){\rule{0em}{1.75em}};%
}

\newcommand\Overline[3][line width=1pt]{%
    \begin{tikzpicture}[remember picture,overlay]
      \draw[#1] (#2.north east) -- (#3.north west);
    \end{tikzpicture}
    }

\title{Grupy i Ciała}
\date{2026}
% \author{Maciej Różewicz}

\begin{document}
\maketitle

\begin{enumerate}
    \item Dla każdego przypadku określ, czy podany zbiór z operacją tworzy
    grupę.
    Jeśli nie — wskaż konkretny aksjomat, który nie jest spełniony.
    \begin{enumerate}
        \item $(\mathbb Z,+)$,
        \item $(\mathbb N,+)$,
        \item $(\mathbb Z_5,+_5)$,
        \item $(\mathbb Z_5,\cdot_5)$,
        \item $(\mathbb Z_5\setminus\{0\},\cdot_5)$,
        \item $(\mathbb Z_6\setminus\{0\},\cdot_6)$.
    \end{enumerate}

    \item Znajdź element odwrotny dla każdego elementu grupy:
    \begin{enumerate}
        \item $(\mathbb Z_7,+_7)$,
        \item $(\mathbb Z_7^*,\cdot_7)$.
    \end{enumerate}

    \item Szyfr Cezara jako grupa - niech:
    \[ G=(\mathbb Z_{26},+_{26}). \]
    Szyfrowanie polega na:
    \[ E_k(x)=x+_{26}k. \]

    Zaszyfruj:
    \[ \text{ALGEBRA} \]
    kluczem:
    $$k=3.$$
    Następnie znajdź operację odszyfrowującą i odszyfruj wiadomość $JHRPHWULD$.

    Pytanie dodatkowe:
    Jakim jednym kluczem można zastąpić kolejno wykonane szyfrowania:
    $E_7$ oraz $E_{12}$?


    \item Dla każdej z poniższych struktur określ, czy jest ciałem. Uzasadnij odpowiedź.
    \begin{enumerate}
        \item $(\mathbb Z_5,+_5,\cdot_5)$
        \item $(\mathbb Z_6,+_6,\cdot_6)$
        \item $(\mathbb Z_7,+_7,\cdot_7)$
    \end{enumerate}

    \item Rozwiąż równanie:
    $$3x+2=5\pmod 7.$$

    \item Pracujemy w ciele:
    $$\mathbb Z_{11}.$$
    Niech:
    $$ a=7,\qquad b=8. $$
    Oblicz:
    $$ a\cdot b\pmod{11}. $$
    Znajdź element odwrotny do \(a=7\), czyli \(7^{-1}\).

    Wykorzystaj go do rozwiązania równania:
    $$ 7x=8\pmod{11}. $$
    Sprawdź, że wynik spełnia równanie.

    Dodatkowe pytanie:
    Dlaczego podobne operacje są możliwe w \(\mathbb Z_{11}\), ale nie możemy ich wykonywać w ten sam sposób w \(\mathbb Z_{12}\)?
\end{enumerate}

\newpage
\textbf{Odpowiedzi:}
\begin{description}
    \item [1:] a)tak, b)nie, c)tak, d)nie, e)tak, f) nie,
    \item [2:] a)$0 \rightarrow 0$, $1 \rightarrow 6$, $2 \rightarrow 5$, $3 \rightarrow 4$, $4 \rightarrow 3$, $5 \rightarrow 2$, $6 \rightarrow 1$,
    b)$1 \rightarrow 1$, $2 \rightarrow 4$, $3 \rightarrow 5$, $4 \rightarrow 2$, $5 \rightarrow 3$, $6 \rightarrow 6$,
    \item [3:] $DOJHEUD$, $GEOMETRIA$
    \item [4:] a)tak, b)nie, c)tak,
    \item [5:] $x=1$,
    \item [6:] $a\cdot b=1$, $7^{-1} = 8$
\end{description}

\end{document}